\IfFileExists{stacks-project.cls}{%
\documentclass{stacks-project}
}{%
\documentclass{amsart}
}

% For dealing with references we use the comment environment
\usepackage{verbatim}
\newenvironment{reference}{\comment}{\endcomment}
%\newenvironment{reference}{}{}
\newenvironment{slogan}{\comment}{\endcomment}
\newenvironment{history}{\comment}{\endcomment}

% For commutative diagrams we use Xy-pic
\usepackage[all]{xy}

% We use 2cell for 2-commutative diagrams.
\xyoption{2cell}
\UseAllTwocells

% We use multicol for the list of chapters between chapters
\usepackage{multicol}

% This is generally recommended for better output
\usepackage{lmodern}
\usepackage[T1]{fontenc}

% For cross-file-references
\usepackage{xr-hyper}

% Package for hypertext links:
\usepackage{hyperref}

% For any local file, say "hello.tex" you want to link to please
% use \externaldocument[hello-]{hello}
\externaldocument[introduction-]{introduction}
\externaldocument[conventions-]{conventions}
\externaldocument[sets-]{sets}
\externaldocument[categories-]{categories}
\externaldocument[topology-]{topology}
\externaldocument[sheaves-]{sheaves}
\externaldocument[sites-]{sites}
\externaldocument[stacks-]{stacks}
\externaldocument[fields-]{fields}
\externaldocument[algebra-]{algebra}
\externaldocument[brauer-]{brauer}
\externaldocument[homology-]{homology}
\externaldocument[derived-]{derived}
\externaldocument[simplicial-]{simplicial}
\externaldocument[more-algebra-]{more-algebra}
\externaldocument[smoothing-]{smoothing}
\externaldocument[modules-]{modules}
\externaldocument[sites-modules-]{sites-modules}
\externaldocument[injectives-]{injectives}
\externaldocument[cohomology-]{cohomology}
\externaldocument[sites-cohomology-]{sites-cohomology}
\externaldocument[dga-]{dga}
\externaldocument[dpa-]{dpa}
\externaldocument[sdga-]{sdga}
\externaldocument[hypercovering-]{hypercovering}
\externaldocument[schemes-]{schemes}
\externaldocument[constructions-]{constructions}
\externaldocument[properties-]{properties}
\externaldocument[morphisms-]{morphisms}
\externaldocument[coherent-]{coherent}
\externaldocument[divisors-]{divisors}
\externaldocument[limits-]{limits}
\externaldocument[varieties-]{varieties}
\externaldocument[topologies-]{topologies}
\externaldocument[descent-]{descent}
\externaldocument[perfect-]{perfect}
\externaldocument[more-morphisms-]{more-morphisms}
\externaldocument[flat-]{flat}
\externaldocument[groupoids-]{groupoids}
\externaldocument[more-groupoids-]{more-groupoids}
\externaldocument[etale-]{etale}
\externaldocument[chow-]{chow}
\externaldocument[intersection-]{intersection}
\externaldocument[pic-]{pic}
\externaldocument[weil-]{weil}
\externaldocument[adequate-]{adequate}
\externaldocument[dualizing-]{dualizing}
\externaldocument[duality-]{duality}
\externaldocument[discriminant-]{discriminant}
\externaldocument[derham-]{derham}
\externaldocument[local-cohomology-]{local-cohomology}
\externaldocument[algebraization-]{algebraization}
\externaldocument[curves-]{curves}
\externaldocument[resolve-]{resolve}
\externaldocument[models-]{models}
\externaldocument[functors-]{functors}
\externaldocument[equiv-]{equiv}
\externaldocument[pione-]{pione}
\externaldocument[etale-cohomology-]{etale-cohomology}
\externaldocument[proetale-]{proetale}
\externaldocument[relative-cycles-]{relative-cycles}
\externaldocument[more-etale-]{more-etale}
\externaldocument[trace-]{trace}
\externaldocument[crystalline-]{crystalline}
\externaldocument[spaces-]{spaces}
\externaldocument[spaces-properties-]{spaces-properties}
\externaldocument[spaces-morphisms-]{spaces-morphisms}
\externaldocument[decent-spaces-]{decent-spaces}
\externaldocument[spaces-cohomology-]{spaces-cohomology}
\externaldocument[spaces-limits-]{spaces-limits}
\externaldocument[spaces-divisors-]{spaces-divisors}
\externaldocument[spaces-over-fields-]{spaces-over-fields}
\externaldocument[spaces-topologies-]{spaces-topologies}
\externaldocument[spaces-descent-]{spaces-descent}
\externaldocument[spaces-perfect-]{spaces-perfect}
\externaldocument[spaces-more-morphisms-]{spaces-more-morphisms}
\externaldocument[spaces-flat-]{spaces-flat}
\externaldocument[spaces-groupoids-]{spaces-groupoids}
\externaldocument[spaces-more-groupoids-]{spaces-more-groupoids}
\externaldocument[bootstrap-]{bootstrap}
\externaldocument[spaces-pushouts-]{spaces-pushouts}
\externaldocument[spaces-chow-]{spaces-chow}
\externaldocument[groupoids-quotients-]{groupoids-quotients}
\externaldocument[spaces-more-cohomology-]{spaces-more-cohomology}
\externaldocument[spaces-simplicial-]{spaces-simplicial}
\externaldocument[spaces-duality-]{spaces-duality}
\externaldocument[formal-spaces-]{formal-spaces}
\externaldocument[restricted-]{restricted}
\externaldocument[spaces-resolve-]{spaces-resolve}
\externaldocument[formal-defos-]{formal-defos}
\externaldocument[defos-]{defos}
\externaldocument[cotangent-]{cotangent}
\externaldocument[examples-defos-]{examples-defos}
\externaldocument[algebraic-]{algebraic}
\externaldocument[examples-stacks-]{examples-stacks}
\externaldocument[stacks-sheaves-]{stacks-sheaves}
\externaldocument[criteria-]{criteria}
\externaldocument[artin-]{artin}
\externaldocument[quot-]{quot}
\externaldocument[stacks-properties-]{stacks-properties}
\externaldocument[stacks-morphisms-]{stacks-morphisms}
\externaldocument[stacks-limits-]{stacks-limits}
\externaldocument[stacks-cohomology-]{stacks-cohomology}
\externaldocument[stacks-perfect-]{stacks-perfect}
\externaldocument[stacks-introduction-]{stacks-introduction}
\externaldocument[stacks-more-morphisms-]{stacks-more-morphisms}
\externaldocument[stacks-geometry-]{stacks-geometry}
\externaldocument[moduli-]{moduli}
\externaldocument[moduli-curves-]{moduli-curves}
\externaldocument[examples-]{examples}
\externaldocument[exercises-]{exercises}
\externaldocument[guide-]{guide}
\externaldocument[desirables-]{desirables}
\externaldocument[coding-]{coding}
\externaldocument[obsolete-]{obsolete}
\externaldocument[fdl-]{fdl}
\externaldocument[index-]{index}

% Theorem environments.
%
\theoremstyle{plain}
\newtheorem{theorem}[subsection]{Theorem}
\newtheorem{proposition}[subsection]{Proposition}
\newtheorem{lemma}[subsection]{Lemma}

\theoremstyle{definition}
\newtheorem{definition}[subsection]{Definition}
\newtheorem{example}[subsection]{Example}
\newtheorem{exercise}[subsection]{Exercise}
\newtheorem{situation}[subsection]{Situation}

\theoremstyle{remark}
\newtheorem{remark}[subsection]{Remark}
\newtheorem{remarks}[subsection]{Remarks}

\numberwithin{equation}{subsection}

% Macros
%
\def\lim{\mathop{\mathrm{lim}}\nolimits}
\def\colim{\mathop{\mathrm{colim}}\nolimits}
\def\Spec{\mathop{\mathrm{Spec}}}
\def\Hom{\mathop{\mathrm{Hom}}\nolimits}
\def\Ext{\mathop{\mathrm{Ext}}\nolimits}
\def\SheafHom{\mathop{\mathcal{H}\!\mathit{om}}\nolimits}
\def\SheafExt{\mathop{\mathcal{E}\!\mathit{xt}}\nolimits}
\def\Sch{\mathit{Sch}}
\def\Mor{\mathop{\mathrm{Mor}}\nolimits}
\def\Ob{\mathop{\mathrm{Ob}}\nolimits}
\def\Sh{\mathop{\mathit{Sh}}\nolimits}
\def\NL{\mathop{N\!L}\nolimits}
\def\CH{\mathop{\mathrm{CH}}\nolimits}
\def\proetale{{pro\text{-}\acute{e}tale}}
\def\etale{{\acute{e}tale}}
\def\QCoh{\mathit{QCoh}}
\def\Ker{\mathop{\mathrm{Ker}}}
\def\Im{\mathop{\mathrm{Im}}}
\def\Coker{\mathop{\mathrm{Coker}}}
\def\Coim{\mathop{\mathrm{Coim}}}

% Boxtimes
%
\DeclareMathSymbol{\boxtimes}{\mathbin}{AMSa}{"02}

%
% Macros for moduli stacks/spaces
%
\def\QCohstack{\mathcal{QC}\!\mathit{oh}}
\def\Cohstack{\mathcal{C}\!\mathit{oh}}
\def\Spacesstack{\mathcal{S}\!\mathit{paces}}
\def\Quotfunctor{\mathrm{Quot}}
\def\Hilbfunctor{\mathrm{Hilb}}
\def\Curvesstack{\mathcal{C}\!\mathit{urves}}
\def\Polarizedstack{\mathcal{P}\!\mathit{olarized}}
\def\Complexesstack{\mathcal{C}\!\mathit{omplexes}}
% \Pic is the operator that assigns to X its picard group, usage \Pic(X)
% \Picardstack_{X/B} denotes the Picard stack of X over B
% \Picardfunctor_{X/B} denotes the Picard functor of X over B
\def\Pic{\mathop{\mathrm{Pic}}\nolimits}
\def\Picardstack{\mathcal{P}\!\mathit{ic}}
\def\Picardfunctor{\mathrm{Pic}}
\def\Deformationcategory{\mathcal{D}\!\mathit{ef}}
\usepackage{cleveref}
\crefname{originallemma}{original lemma}{original lemmas}
\Crefname{originallemma}{Original lemma}{Original lemmas}
\makeatletter
\newlabel{lemma-finite-stratify}{{0AB9}{source}{}{sourcepointer-1}{}}
\newlabel{lemma-finite-stratify@cref}{{[originallemma][0][]0AB9}{[1][0][]source}{}{}{}}
\makeatother
\title{R-invariant thickening decomposition of U for a groupoid scheme with finite source and target, with finite locally free restriction on the closed part}
\author{The Stacks Project Authors}
\date{stacks-project:more-groupoids.tex:local-snapshot-20261012; GNU FDL 1.2 or later}
\begin{document}
\maketitle
\setcounter{subsection}{1}
\noindent\textit{Editorial note.} All delivered mathematical text is the author's own, copied from the pinned Stacks Project chapter file more-groupoids.tex; no statement or proof step was rewritten, reordered or removed. The counted span is the author environment and its own proof environment. References inside the span are handed to the assembler as navigation stubs that carry the canonical tag of the cited tag only; whether each of them is load-bearing is recorded on the stub itself.\par\medskip
\noindent\textit{Allowed standard prerequisites (editorial).}\begin{itemize}
\item The finite stratification of U for a groupoid scheme with finite source and target: U = Z\_0 superset Z\_1 superset Z\_2 superset ... R-invariant closed subschemes with empty intersection, whose strata differences carry restrictions with finite locally free source and target of rank r (more-groupoids chapter, Lemma lemma-finite-stratify, canonical tag 0AB9). cited at proof lines 2686-2687, 2735-2736 and 2739. Supplies the stratification that drives the construction, the rank r finite locally free property (5) on the new stratum, and the empty intersection that yields the covering property (2).
\end{itemize}\medskip
\ifdefined\hypertarget\hypertarget{sourcepointer-1}{}\fi
\noindent\textit{Source reference.} Original lemma 0AB9: sources/stacks-chow/tags.txt, source lines 13012--13012. Cited at proof lines 2686-2687 for the stratification U = Z\_0 superset Z\_1 superset ... that drives the induction, at lines 2735-2736 for the rank r finite locally free property (5) on the new stratum, and at line 2739 for the empty intersection of the strata giving the covering (2).\par\medskip
\begin{lemma}
\label{lemma-finite-flat-over-almost-dense-subscheme}
Let $(U, R, s, t, c)$ be a groupoid scheme over a scheme $S$. Assume $s, t$
are finite. There exists an open subscheme $W \subset U$ and a closed
subscheme $W' \subset W$ such that
\begin{enumerate}
\item $W$ and $W'$ are $R$-invariant,
\item $U = t(s^{-1}(\overline{W}))$ set theoretically,
\item $W$ is a thickening of $W'$, and
\item the maps $s'$, $t'$ of the restriction $(W', R', s', t', c')$
are finite locally free.
\end{enumerate}
\end{lemma}

\begin{proof}
Consider the stratification $U = Z_0 \supset Z_1 \supset Z_2 \supset \ldots$
of Lemma \ref{lemma-finite-stratify}.

\medskip\noindent
We will construct disjoint unions $W = \coprod_{r \geq 1} W_r$ and
$W' = \coprod_{r \geq 1} W'_r$ with each $W'_r \to W_r$ a thickening
of $R$-invariant subschemes of $U$ such that the morphisms
$s_r', t_r'$ of the restrictions $(W_r', R_r', s_r', t_r', c_r')$
are finite locally free of rank $r$. To begin we set
$W_1 = W'_1 = U \setminus Z_1$. This is an $R$-invariant open
subscheme of $U$, it is true that $W_0$ is a thickening of $W'_0$,
and the maps $s_1'$, $t_1'$ of the
restriction $(W_1', R_1', s_1', t_1', c_1')$ are isomorphisms, i.e.,
finite locally free of rank $1$.
Moreover, every point of $U \setminus Z_1$ is in $t(s^{-1}(\overline{W_1}))$.

\medskip\noindent
Assume we have found subschemes $W'_r \subset W_r \subset U$ for $r \leq n$
such that
\begin{enumerate}
\item $W_1, \ldots, W_n$ are disjoint,
\item $W_r$ and $W_r'$ are $R$-invariant,
\item $U \setminus Z_n \subset \bigcup_{r \leq n} t(s^{-1}(\overline{W_r}))$
set theoretically,
\item $W_r$ is a thickening of $W'_r$,
\item the maps $s_r'$, $t_r'$ of the restriction
$(W_r', R_r', s_r', t_r', c_r')$ are finite locally free of rank $r$.
\end{enumerate}
Then we set
$$
W_{n + 1} = Z_n \setminus
\left(
Z_{n + 1} \cup \bigcup\nolimits_{r \leq n} t(s^{-1}(\overline{W_r}))
\right)
$$
set theoretically and
$$
W'_{n + 1} = Z_n \setminus
\left(
Z_{n + 1} \cup \bigcup\nolimits_{r \leq n} t(s^{-1}(\overline{W_r}))
\right)
$$
scheme theoretically. Then $W_{n + 1}$ is an $R$-invariant open subscheme
of $U$ because $Z_{n + 1} \setminus \overline{U \setminus Z_{n + 1}}$
is open in $U$ and $\overline{U \setminus Z_{n + 1}}$ is contained
in the closed subset $\bigcup\nolimits_{r \leq n} t(s^{-1}(\overline{W_r}))$
we are removing by property (3) and the fact that $t$ is a closed morphism.
It is clear that $W'_{n + 1}$ is a closed subscheme
of $W_{n + 1}$ with the same underlying topological space.
Finally, properties (1), (2) and (3) are clear and property (5) follows from
Lemma \ref{lemma-finite-stratify}.

\medskip\noindent
By Lemma \ref{lemma-finite-stratify} we have $\bigcap Z_r = \emptyset$.
Hence every point of $U$ is contained in $U \setminus Z_n$
for some $n$. Thus we see that
$U = \bigcup_{r \geq 1} t(s^{-1}(\overline{W_r}))$
set theoretically and we see that (2) holds.
Thus $W' \subset W$ satisfy (1), (2), (3), and (4).
\end{proof}

\end{document}
